\chapter{Communication Complexity — Finite Message}
%%% generated by the blueprint pipeline from TCSlib.CommunicationComplexity.DeterministicCC.FiniteMessage
\section{Overview}

This module develops a generalized model of deterministic two-party communication protocols in
which, at each step, a player sends an element of an arbitrary finite type $\beta$ rather than
a single bit.  The central results establish that every such finite-message protocol can be
converted to an equivalent binary protocol with the same run behavior and the same communication
complexity, where each $\beta$-valued message is encoded using $\lceil \log_2 |\beta| \rceil$
bits.

%%%%%%%%%%%%%%%%%%%%%%%%%%%%%%%%%%%%%%%%%%%%%%%%%%%%%%%%%%%%%%%%%%%%%%%%%%%%%%%
\section{Declarations}

\begin{definition}[Finite-message protocol]
\lean{CommunicationComplexity.Deterministic.FiniteMessage.Protocol}
\label{CommunicationComplexity.Deterministic.FiniteMessage.Protocol}
\leanok
A generalized deterministic two-party communication protocol over input types $X$, $Y$ and
output type $\alpha$.  At each step a player sends an element of an arbitrary finite nonempty
type $\beta$: the \texttt{alice} constructor takes a function $f : X \to \beta$ and a
continuation $P : \beta \to \mathrm{Protocol}\,X\,Y\,\alpha$, and dually for \texttt{bob}.
The \texttt{output} constructor terminates the protocol with a value in $\alpha$.  This
inductive type is equivalent to the binary \texttt{Deterministic.Protocol} up to complexity,
where a $\beta$-valued message costs $\lceil \log_2 |\beta| \rceil$ bits.
\end{definition}

\begin{definition}[Protocol execution]
\lean{CommunicationComplexity.Deterministic.FiniteMessage.Protocol.run}
\label{CommunicationComplexity.Deterministic.FiniteMessage.Protocol.run}
\leanok
\uses{CommunicationComplexity.Deterministic.FiniteMessage.Protocol}
Given a finite-message protocol $p$ and inputs $x : X$, $y : Y$, \texttt{run} executes $p$
recursively: at an \texttt{alice} node it evaluates $f(x)$ and follows the corresponding
continuation, at a \texttt{bob} node it evaluates $f(y)$, and at an \texttt{output} node it
returns the stored value.
\end{definition}

\begin{definition}[Communication complexity of a finite-message protocol]
\lean{CommunicationComplexity.Deterministic.FiniteMessage.Protocol.complexity}
\label{CommunicationComplexity.Deterministic.FiniteMessage.Protocol.complexity}
\leanok
\uses{CommunicationComplexity.Deterministic.FiniteMessage.Protocol}
The worst-case communication cost of a finite-message protocol, measured in bits.  An
\texttt{output} node costs $0$; an \texttt{alice} or \texttt{bob} node with message type
$\beta$ costs $\lceil \log_2 |\beta| \rceil$ plus the maximum complexity over all continuations,
i.e.\ $\lceil \log_2 |\beta| \rceil + \sup_{b \in \beta} \mathrm{complexity}(P\,b)$.
\end{definition}

\begin{definition}[Complete binary tree for Alice queries]
\lean{CommunicationComplexity.Deterministic.FiniteMessage.Protocol.completeTreeAlice}
\label{CommunicationComplexity.Deterministic.FiniteMessage.Protocol.completeTreeAlice}
\leanok
\uses{CommunicationComplexity.Deterministic.Protocol}
An auxiliary construction: given $d$ binary query functions $\mathrm{query}_i : X \to
\mathrm{Bool}$ and a family of binary protocols $Q : (\mathrm{Fin}\,d \to \mathrm{Bool}) \to
\mathrm{Protocol}\,X\,Y\,\alpha$, builds a single binary protocol that reads each query bit
from Alice in sequence and then runs $Q$ on the collected bit pattern.
\end{definition}

\begin{theorem}[Evaluation of the complete Alice-query tree]
\lean{CommunicationComplexity.Deterministic.FiniteMessage.Protocol.completeTreeAlice_run}
\label{CommunicationComplexity.Deterministic.FiniteMessage.Protocol.completeTreeAlice_run}
\leanok
\uses{CommunicationComplexity.Deterministic.FiniteMessage.Protocol.completeTreeAlice, CommunicationComplexity.Deterministic.Protocol, CommunicationComplexity.Deterministic.Protocol.run}
\difficulty{4}
Fix a depth $d \in \bbn$, a family of binary query functions $\mathrm{query}_i : X \to
\mathrm{Bool}$ for $i \in \mathrm{Fin}\,d$, and a family of deterministic protocols
$Q(b)$ indexed by bit patterns $b : \mathrm{Fin}\,d \to \mathrm{Bool}$. Then for every
pair of inputs $x : X$ and $y : Y$, the protocol that first has Alice read the $d$ query
bits in sequence and then runs $Q$ on the collected pattern evaluates on $(x,y)$ to the
same value as the single protocol $Q\bigl(i \mapsto \mathrm{query}_i(x)\bigr)$ evaluated
on $(x,y)$.
\end{theorem}

\begin{theorem}[Complexity of the complete Alice-query tree]
\lean{CommunicationComplexity.Deterministic.FiniteMessage.Protocol.completeTreeAlice_complexity}
\label{CommunicationComplexity.Deterministic.FiniteMessage.Protocol.completeTreeAlice_complexity}
\leanok
\uses{CommunicationComplexity.Deterministic.FiniteMessage.Protocol.completeTreeAlice, CommunicationComplexity.Deterministic.Protocol, CommunicationComplexity.Deterministic.Protocol.complexity}
\difficulty{5}
Fix a natural number $d$, together with binary query functions
$\mathrm{query}_0,\dots,\mathrm{query}_{d-1} : X \to \mathrm{Bool}$ and a family of
deterministic protocols $Q(\mathrm{bits})$ indexed by the bit patterns $\mathrm{bits} :
\mathrm{Fin}\,d \to \mathrm{Bool}$. Form the protocol that has Alice read the bits
$\mathrm{query}_0(x),\dots,\mathrm{query}_{d-1}(x)$ in sequence and then run $Q$ on the
collected pattern. Its communication complexity is
\[
d + \max_{\mathrm{bits} : \mathrm{Fin}\,d \to \mathrm{Bool}} \;
\mathrm{complexity}\bigl(Q(\mathrm{bits})\bigr),
\]
the maximum being taken over all $2^d$ length-$d$ bit patterns.
\end{theorem}

\begin{theorem}[Binary encoding of a finite-alphabet message from Alice]
\lean{CommunicationComplexity.Deterministic.FiniteMessage.Protocol.encode_alice}
\label{CommunicationComplexity.Deterministic.FiniteMessage.Protocol.encode_alice}
\leanok
\uses{CommunicationComplexity.Deterministic.FiniteMessage.Protocol.completeTreeAlice, CommunicationComplexity.Deterministic.FiniteMessage.Protocol.completeTreeAlice_complexity, CommunicationComplexity.Deterministic.FiniteMessage.Protocol.completeTreeAlice_run, CommunicationComplexity.Deterministic.Protocol, CommunicationComplexity.Deterministic.Protocol.complexity, CommunicationComplexity.Deterministic.Protocol.run}
\difficulty{6}
Let $X$, $Y$, and $\alpha$ be types, let $\beta$ be a finite nonempty type, let $f : X
\to \beta$ be a function, and let $Q : \beta \to \mathrm{Protocol}\,X\,Y\,\alpha$ assign
to each value $b \in \beta$ a deterministic two-party communication protocol. Then there
exists a deterministic protocol $R$ over the same input types such that, on every pair
of inputs $x : X$ and $y : Y$, running $R$ yields the same output as running the
protocol $Q(f(x))$, and the communication complexity of $R$ equals
\[
\lceil \log_2 \abs{\beta} \rceil \;+\; \max_{b \in \beta}
\operatorname{complexity}(Q(b)),
\]
where $\abs{\beta}$ denotes the cardinality of $\beta$.
\end{theorem}

\begin{theorem}[Binary simulation of a finite-message protocol]
\lean{CommunicationComplexity.Deterministic.FiniteMessage.Protocol.toProtocol_exists}
\label{CommunicationComplexity.Deterministic.FiniteMessage.Protocol.toProtocol_exists}
\leanok
\uses{CommunicationComplexity.Deterministic.FiniteMessage.Protocol, CommunicationComplexity.Deterministic.FiniteMessage.Protocol.complexity, CommunicationComplexity.Deterministic.FiniteMessage.Protocol.encode_alice, CommunicationComplexity.Deterministic.FiniteMessage.Protocol.run, CommunicationComplexity.Deterministic.Protocol, CommunicationComplexity.Deterministic.Protocol.complexity, CommunicationComplexity.Deterministic.Protocol.run, CommunicationComplexity.Deterministic.Protocol.swap, CommunicationComplexity.Deterministic.Protocol.swap_complexity, CommunicationComplexity.Deterministic.Protocol.swap_run}
\difficulty{4}
Let $X$, $Y$, and $\alpha$ be types, and let $p$ be a finite-message two-party protocol
over inputs $X$ (Alice) and $Y$ (Bob) with output type $\alpha$, in which each message
ranges over an arbitrary finite nonempty type and a $\beta$-valued message is charged
$\lceil \log_2 |\beta| \rceil$ bits. Then there exists a binary deterministic protocol
$P$ over the same input types $X$, $Y$ and output type $\alpha$ — one in which every
message is a single bit — that computes the same function, i.e.\ $P$ and $p$ return the
same output on every pair of inputs, and has exactly the same communication complexity
as $p$.
\end{theorem}

\begin{definition}[Conversion to binary protocol]
\lean{CommunicationComplexity.Deterministic.FiniteMessage.Protocol.toProtocol}
\label{CommunicationComplexity.Deterministic.FiniteMessage.Protocol.toProtocol}
\leanok
\uses{CommunicationComplexity.Deterministic.FiniteMessage.Protocol, CommunicationComplexity.Deterministic.FiniteMessage.Protocol.toProtocol_exists, CommunicationComplexity.Deterministic.Protocol}
A noncomputable function that converts a finite-message protocol $p$ into a binary protocol
\texttt{CommunicationComplexity.Deterministic.FiniteMessage.Protocol.toProtocol} $p$ with the same run behavior and the same communication complexity,
encoding each $\beta$-valued message as $\lceil \log_2 |\beta| \rceil$ bits.
\end{definition}

\begin{theorem}[Conversion to binary protocol preserves output]
\lean{CommunicationComplexity.Deterministic.FiniteMessage.Protocol.toProtocol_run}
\label{CommunicationComplexity.Deterministic.FiniteMessage.Protocol.toProtocol_run}
\leanok
\uses{CommunicationComplexity.Deterministic.FiniteMessage.Protocol, CommunicationComplexity.Deterministic.FiniteMessage.Protocol.run, CommunicationComplexity.Deterministic.FiniteMessage.Protocol.toProtocol, CommunicationComplexity.Deterministic.FiniteMessage.Protocol.toProtocol_exists, CommunicationComplexity.Deterministic.Protocol.run}
\difficulty{1}
Let $p$ be a finite-message deterministic two-party protocol over input types $X$ and
$Y$ with output type $\alpha$, and let $q$ be the binary protocol obtained by converting
$p$ (encoding each message drawn from a finite nonempty type $\beta$ as $\lceil \log_2
\abs{\beta} \rceil$ bits). Then $q$ computes the same function as $p$: for every pair of
inputs $x \in X$ and $y \in Y$, running $q$ on $(x, y)$ returns the same value in
$\alpha$ as running $p$ on $(x, y)$.
\end{theorem}

\begin{theorem}[Conversion to a binary protocol preserves complexity]
\lean{CommunicationComplexity.Deterministic.FiniteMessage.Protocol.toProtocol_complexity}
\label{CommunicationComplexity.Deterministic.FiniteMessage.Protocol.toProtocol_complexity}
\leanok
\uses{CommunicationComplexity.Deterministic.FiniteMessage.Protocol, CommunicationComplexity.Deterministic.FiniteMessage.Protocol.complexity, CommunicationComplexity.Deterministic.FiniteMessage.Protocol.toProtocol, CommunicationComplexity.Deterministic.FiniteMessage.Protocol.toProtocol_exists, CommunicationComplexity.Deterministic.Protocol.complexity}
\difficulty{1}
Let $p$ be a finite-message deterministic two-party communication protocol over input
types $X$ (Alice) and $Y$ (Bob) with outputs in $\alpha$, in which each message is an
element of an arbitrary finite nonempty type. Its conversion into a binary protocol,
obtained by encoding every $\beta$-valued message as $\lceil \log_2 \abs{\beta} \rceil$
bits, has communication complexity equal to that of $p$.
\end{theorem}

\begin{definition}[Embedding binary protocols as finite-message protocols]
\lean{CommunicationComplexity.Deterministic.FiniteMessage.Protocol.ofProtocol}
\label{CommunicationComplexity.Deterministic.FiniteMessage.Protocol.ofProtocol}
\leanok
\uses{CommunicationComplexity.Deterministic.FiniteMessage.Protocol, CommunicationComplexity.Deterministic.Protocol}
Embeds a binary protocol into the generalized finite-message framework by treating each
Boolean message as an element of $\beta = \mathrm{Bool}$, giving a
\texttt{CommunicationComplexity.Deterministic.Protocol} $X\,Y\,\alpha$ with the same tree structure.
\end{definition}

\begin{theorem}[Embedding preserves protocol output]
\lean{CommunicationComplexity.Deterministic.FiniteMessage.Protocol.ofProtocol_run}
\label{CommunicationComplexity.Deterministic.FiniteMessage.Protocol.ofProtocol_run}
\leanok
\uses{CommunicationComplexity.Deterministic.FiniteMessage.Protocol.ofProtocol, CommunicationComplexity.Deterministic.FiniteMessage.Protocol.run, CommunicationComplexity.Deterministic.Protocol, CommunicationComplexity.Deterministic.Protocol.run}
\difficulty{2}
Let $X$ and $Y$ be the input types of Alice and Bob, let $\alpha$ be an output type, and
let $p$ be a deterministic binary two-party communication protocol producing values in
$\alpha$. Regard $p$ as a finite-message protocol by treating each of its Boolean
messages as an element of the two-element type. Then for all inputs $x \in X$ and $y \in
Y$, executing this embedded finite-message protocol on $x$ and $y$ yields the same value
in $\alpha$ as executing $p$ directly on $x$ and $y$.
\end{theorem}

\begin{theorem}[Embedding binary protocols preserves complexity]
\lean{CommunicationComplexity.Deterministic.FiniteMessage.Protocol.ofProtocol_complexity}
\label{CommunicationComplexity.Deterministic.FiniteMessage.Protocol.ofProtocol_complexity}
\leanok
\uses{CommunicationComplexity.Deterministic.FiniteMessage.Protocol.complexity, CommunicationComplexity.Deterministic.FiniteMessage.Protocol.ofProtocol, CommunicationComplexity.Deterministic.Protocol, CommunicationComplexity.Deterministic.Protocol.complexity}
\difficulty{2}
Let $p$ be a deterministic binary two-party communication protocol over input types $X$
and $Y$ with output type $\alpha$, and let its embedding into the finite-message
framework send each Boolean message as an element of $\{\mathtt{false},\mathtt{true}\}$.
Then the communication complexity of the embedded finite-message protocol equals the
communication complexity of $p$.
\end{theorem}

\begin{theorem}[Realizing binary protocols as finite-message protocols]
\lean{CommunicationComplexity.Deterministic.FiniteMessage.Protocol.ofProtocol_equiv}
\label{CommunicationComplexity.Deterministic.FiniteMessage.Protocol.ofProtocol_equiv}
\leanok
\uses{CommunicationComplexity.Deterministic.FiniteMessage.Protocol, CommunicationComplexity.Deterministic.FiniteMessage.Protocol.complexity, CommunicationComplexity.Deterministic.FiniteMessage.Protocol.ofProtocol, CommunicationComplexity.Deterministic.FiniteMessage.Protocol.ofProtocol_complexity, CommunicationComplexity.Deterministic.FiniteMessage.Protocol.ofProtocol_run, CommunicationComplexity.Deterministic.FiniteMessage.Protocol.run, CommunicationComplexity.Deterministic.Protocol, CommunicationComplexity.Deterministic.Protocol.complexity, CommunicationComplexity.Deterministic.Protocol.run}
\difficulty{1}
Let $X$ and $Y$ be the input types of Alice and Bob and let $\alpha$ be an output type.
For every binary deterministic communication protocol $p$ over $X$, $Y$, and $\alpha$,
there is a finite-message protocol $P$ over the same input and output types that
computes the same output on every pair of inputs — that is, $\mathrm{run}(P) =
\mathrm{run}(p)$ as functions of $x \in X$ and $y \in Y$ — and has the same
communication complexity, $\mathrm{complexity}(P) = \mathrm{complexity}(p)$.
\end{theorem}

\begin{definition}[Pullback along input maps]
\lean{CommunicationComplexity.Deterministic.FiniteMessage.Protocol.comap}
\label{CommunicationComplexity.Deterministic.FiniteMessage.Protocol.comap}
\leanok
\uses{CommunicationComplexity.Deterministic.FiniteMessage.Protocol}
Given maps $f_X : X' \to X$ and $f_Y : Y' \to Y$, the pullback
$p.\mathrm{comap}\,f_X\,f_Y$ is the finite-message protocol over $X'$, $Y'$ obtained by
precomposing every message function with $f_X$ or $f_Y$ respectively, leaving the protocol
tree structure unchanged.
\end{definition}

\begin{theorem}[Execution commutes with pullback]
\lean{CommunicationComplexity.Deterministic.FiniteMessage.Protocol.comap_run}
\label{CommunicationComplexity.Deterministic.FiniteMessage.Protocol.comap_run}
\leanok
\uses{CommunicationComplexity.Deterministic.FiniteMessage.Protocol, CommunicationComplexity.Deterministic.FiniteMessage.Protocol.comap, CommunicationComplexity.Deterministic.FiniteMessage.Protocol.run}
\difficulty{2}
Let $p$ be a finite-message deterministic communication protocol over input types $X$,
$Y$ with output type $\alpha$, and let $f_X : X' \to X$ and $f_Y : Y' \to Y$ be maps
into those input types. Then for every pair of inputs $x' : X'$ and $y' : Y'$, executing
the pullback protocol $p.\mathrm{comap}\,f_X\,f_Y$ on $(x', y')$ yields the same output
as executing $p$ on the transported inputs $(f_X(x'), f_Y(y'))$; that is,
\[
\mathrm{run}\bigl(p.\mathrm{comap}\,f_X\,f_Y\bigr)(x', y') \;=\;
\mathrm{run}(p)\bigl(f_X(x'),\, f_Y(y')\bigr).
\]
\end{theorem}

\begin{theorem}[Pullback preserves communication complexity]
\lean{CommunicationComplexity.Deterministic.FiniteMessage.Protocol.comap_complexity}
\label{CommunicationComplexity.Deterministic.FiniteMessage.Protocol.comap_complexity}
\leanok
\uses{CommunicationComplexity.Deterministic.FiniteMessage.Protocol, CommunicationComplexity.Deterministic.FiniteMessage.Protocol.comap, CommunicationComplexity.Deterministic.FiniteMessage.Protocol.complexity}
\difficulty{2}
Let $p$ be a finite-message deterministic communication protocol over input types $X$
and $Y$ with output type $\alpha$, and let $f_X : X' \to X$ and $f_Y : Y' \to Y$ be maps
of input types. Then the pullback protocol obtained by precomposing every message
function of $p$ with $f_X$ or $f_Y$ (as appropriate) has the same communication
complexity as $p$.
\end{theorem}
